\documentclass[aspectratio=169,10pt]{beamer}
\usetheme[numbering=fraction,progressbar=frametitle]{metropolis}
\usepackage{fontspec}
\setsansfont{Latin Modern Sans}
\usepackage{amsmath,amssymb,mathtools,graphicx,booktabs}
\usepackage{polyglossia}
\setdefaultlanguage{english}
\setbeamertemplate{navigation symbols}{}
\setbeamerfont{frametitle}{size=\large}
\setbeamerfont{normal text}{size=\normalsize}
\setlength{\parskip}{.45em}
\hypersetup{pdfauthor={},pdfsubject={Lecture slides},pdfkeywords={},pdftitle={Synthesis, reasoning and fluency}}
\providecommand{\N}{\mathbb{N}}
\providecommand{\Z}{\mathbb{Z}}
\providecommand{\D}{\mathbb{D}}
\providecommand{\Q}{\mathbb{Q}}
\providecommand{\R}{\mathbb{R}}
\providecommand{\C}{\mathbb{C}}
\providecommand{\vect}[1]{\overrightarrow{#1}}
\providecommand{\vi}{\vec{\imath}}
\providecommand{\vj}{\vec{\jmath}}
\providecommand{\vk}{\vec{k}}
\providecommand{\norm}[1]{\left\lVert #1\right\rVert}
\providecommand{\vc}[2]{\begin{pmatrix}#1\\#2\end{pmatrix}}
\providecommand{\deter}[2]{\det\left(#1,#2\right)}
\providecommand{\abs}[1]{\left\lvert #1\right\rvert}
\providecommand{\dd}{\,\mathrm{d}}
\providecommand{\e}{\mathrm{e}}
\providecommand{\Card}{\operatorname{Card}}
\providecommand{\E}{\operatorname{E}}
\providecommand{\V}{\operatorname{V}}
\providecommand{\degre}{^{\circ}}
\providecommand{\un}[1]{\,\mathrm{#1}}
\providecommand{\sci}[2]{#1\times10^{#2}}
\providecommand{\pgcd}{\operatorname{PGCD}}
\providecommand{\ppcm}{\operatorname{PPCM}}
\title{Synthesis, reasoning and fluency}
\subtitle{Lecture and worked examples}
\date{}
\author{}
\begin{document}
\begin{frame}[plain]
\vfill
{\usebeamerfont{title}\inserttitle\par}
\medskip
{\usebeamerfont{subtitle}\insertsubtitle\par}
\medskip
{\usebeamercolor[fg]{progress bar}\rule{\textwidth}{.5pt}}
\vfill
\end{frame}
\begin{frame}{Implication, converse and counterexample}
“If \(P\), then \(Q\)” is an implication \(P\Rightarrow Q\). Its converse is \(Q\Rightarrow P\); it is not automatically true. An equivalence \(P\iff Q\) requires both implications.
A counterexample is enough to refute a universal statement; a few examples are not enough to prove it. The contrapositive \(\text{not }Q\Rightarrow\text{not }P\) is equivalent to \(P\Rightarrow Q\).
\end{frame}
\begin{frame}{Implication, converse and counterexample}
For an integer \(n\): “If \(n\) is a multiple of \(4\), then \(n\) is even.” Prove the implication, write its converse and its contrapositive, then say which are true.
\end{frame}
\begin{frame}{Worked answer: Implication, converse and counterexample}
If \(n=4k\), then \(n=2(2k)\) is even. Converse: “If \(n\) is even, then \(n\) is a multiple of \(4\)”; it is false with \(n=2\). Contrapositive: “If \(n\) is odd, then \(n\) is not a multiple of \(4\)”; it is true and equivalent to the original implication.
\end{frame}
\begin{frame}{Implication, converse and counterexample}
For \(x\in\R\), study “\(x=3\Rightarrow x^2=9\)” and its converse. Give a correct equivalence.
\end{frame}
\begin{frame}{Worked answer: Implication, converse and counterexample}
The implication is true. The converse is false with \(x=-3\). The correct equivalence is \(x^2=9\iff x=3\text{ or }x=-3\). With the extra hypothesis \(x\geq0\), one can write \(x^2=9\iff x=3\).
\end{frame}
\begin{frame}{Negation, case analysis and proof}
Negate correctly: “\(x>2\)”; “\(x\in[1,4]\)”; “All students play sport”; “\(A\) and \(B\) both occur”.
\end{frame}
\begin{frame}{Worked answer: Negation, case analysis and proof}
“\(x\leq2\)”; “\(x<1\) or \(x>4\)”; “At least one student does not play sport”; “\(A\) does not occur or \(B\) does not occur”. The “or” is inclusive.
\end{frame}
\begin{frame}{Negation, case analysis and proof}
Show by case analysis that \(n(n+1)\) is even for every integer \(n\). Then show that \(n^2+n+1\) is odd.
\end{frame}
\begin{frame}{Worked answer: Negation, case analysis and proof}
If \(n\) is even, \(n(n+1)\) has an even factor. If \(n\) is odd, \(n+1\) is even; the product is again even. The two cases cover all integers. Since \(n^2+n+1=n(n+1)+1\), it is an odd integer.
\end{frame}
\begin{frame}{Negation, case analysis and proof}
A student writes “\(n^2\) even implies \(n\) even, because I checked it for \(n=2,4,6\)”. Replace this justification with a proof.
\end{frame}
\begin{frame}{Worked answer: Negation, case analysis and proof}
If \(n\) were odd, the result proved in sequence 1 would give \(n^2\) odd, a contradiction. This is a contrapositive argument, or presented here as a proof by contradiction. Three examples do not cover all integers.
\end{frame}
\begin{frame}{Combining algebra, functions and percentages}
Two fictional rental rates are offered for \(x\) hours, with \(x\in[0,8]\): \(A(x)=3000+2000x\) Ar and \(B(x)=3500x\) Ar.

\begin{enumerate}
\setcounter{enumi}{0}\item Calculate the costs for \(1\) h and \(4\) h.
\item Find where the rates are equal, then the durations for which \(A\) is cheaper.
\item A \(10\%\) discount is granted on the whole of rate \(A\). Give the new rate \(A_r\) and the new comparison threshold.

\end{enumerate}
\end{frame}
\begin{frame}{Worked answer: Combining algebra, functions and percentages}
1. At \(1\) h: \(5000\) Ar versus \(3500\) Ar. At \(4\) h: \(11000\) Ar versus \(14000\) Ar. 2. \(3000+2000x=3500x\iff x=2\). \(A<B\) for \(x>2\), so on \((2,8]\). 3. \(A_r(x)=0.9(3000+2000x)=2700+1800x\). Equality with \(B\) gives \(1700x=2700\), that is \(x=27/17\approx1.59\) h. \(A_r<B\) on \((27/17,8]\). The discount also applies to the fixed charge.
\end{frame}
\begin{frame}[fragile]{Python}
\small
\begin{verbatim}
def tarif_moins_cher(x):
    a = 3000 + 2000*x
    b = 3500*x
    if a < b:
        return "A"
    elif b < a:
        return "B"
    else:
        return "egalite"
\end{verbatim}
\end{frame}
\begin{frame}{Combining algebra, functions and percentages}
Predict the results for \(x=1,2,4\). Modify the line that computes \(a\) to take the discount into account.
\end{frame}
\begin{frame}{Worked answer: Combining algebra, functions and percentages}
Results: \texttt{"B"}, \texttt{"egalite"}, \texttt{"A"}. Replace with \texttt{a = 0.9*(3000 + 2000*x)}. The domain \([0,8]\) is assumed to be respected in the calls.
\end{frame}
\end{document}
